% Claim:
%   \texttt{Session} is the first-class runtime value that flows through agents,
%   tools, evidence-gated memory, sandboxes, compressors, and workflows.
% Evidence:
%   initial-tech-report/references/OpenRath-README.md lines 22--84 and 170--241.
% Open questions:
%   Later architecture sections should pin the invariants for merge, compression,
%   replay, and evidence export.

The central claim is that agent systems benefit from a first-class runtime state, and OpenRath proposes \texttt{Session} as that state. To clarify the role of this object, OpenRath adopts a PyTorch-inspired programming model~\cite{pytorch}: not PyTorch's tensor mathematics, but its architectural interface for composable computation. In that interface, a central value flows through reusable modules, modules expose a uniform \texttt{forward} mapping, placement is made explicit through operations such as \texttt{tensor.to(device)}, and persistent module state is represented by parameters. OpenRath adapts this pattern to agent runtimes. \texttt{Session} plays the role of the flowing value; \texttt{Agent} is a reusable transformation similar in role to a layer; \texttt{Workflow} is a compositional container; both follow a \texttt{forward(session) -> session} contract; placement is expressed as \texttt{session.to(backend)}; and \texttt{Memory} is treated as an agent-bound persistent state plane rather than hidden prompt text.

Because each transformation preserves the \texttt{Session -> Session} shape, agents can be nested into workflows without introducing a separate runtime state format. Composition, branching, merging, handoff, and replay operate on ordinary program values rather than on state reconstructed from controller logs. The analogy is architectural rather than literal: the claim is not that agent systems are neural networks, but that agent runtimes need a stable flowing value, reusable transformations behind a uniform interface, explicit placement, persistent state, and inspectable evidence. The compact vocabulary that realizes this design---\texttt{Agent}, \texttt{Workflow}, \texttt{Tool}, \texttt{Memory}, and \texttt{Sandbox}---is developed in the programming model that follows.

Why make \texttt{Session} the runtime boundary, rather than placing this state inside a graph runtime's node state or a tracing system's spans? These layers serve different primary readers. Graph state records where execution is in a control flow, so that a run can resume, replay, or fork from checkpoints. Trace spans record what was observed during execution, such as model calls, tool calls, handoffs, guardrails, and other monitored events. Neither representation is designed to be the ordinary program value that agents themselves fork, merge, hand off, and replay. A trace is written for observers; a graph checkpoint is written for schedulers. A \texttt{Session} is written for the agent program: it is the live value passed through the program, and evidence is attached to that value rather than reconstructed from a side channel. OpenRath's design hypothesis is that multi-agent systems remain more inspectable as they scale when runtime state is placed where the program already flows, rather than beside it.

\begin{table}[!htbp] \vspace{0.1em} \centering \small \caption{Three runtime records, three readers. OpenRath's \texttt{Session} is the value written for the agent program itself, which is why fork, merge, and replay are first-class rather than reconstructed.} \label{tab:session-vs-others} \begin{tabularx}{\linewidth}{@{}>{\raggedright\arraybackslash}p{0.22\linewidth}>{\raggedright\arraybackslash}X>{\raggedright\arraybackslash}X@{}} \toprule Record & Written for & What it primarily holds \\ \midrule Graph checkpoint & The scheduler & Where execution is in the control flow, so a run can resume or time-travel. \\ Trace span & The observer & What was observed during a run, for monitoring and debugging after the fact. \\ \texttt{Session} (OpenRath) & The agent program & The live value agents fork, merge, hand off, and replay; lineage, tools, placement, memory, and usage travel with it. \\ \bottomrule \end{tabularx} \vspace{0.1em} \end{table}